\documentclass[11pt,a4paper]{article}
\usepackage[utf8]{inputenc}
\usepackage[T1]{fontenc}
\usepackage[english]{babel}
\usepackage{geometry}
\geometry{margin=2cm}
\usepackage{graphicx}
\usepackage{hyperref}
\usepackage{xcolor}
\usepackage{amsmath}
\usepackage{booktabs}
\usepackage{longtable}
\usepackage{fancyhdr}

\pagestyle{fancy}
\fancyhf{}
\rhead{Demo Script — Quantum Testing Agent}
\lhead{\thepage}

\title{\textbf{Guion de demo}\\
\large Quantum Testing Agent — Vultr: Agent Rush Hackathon}
\author{Jesús León Porras}
\date{8 de noviembre de 2026}

\begin{document}

\maketitle

\begin{abstract}
Guion de demostración para el proyecto Quantum Testing Agent, participante en el
desafío Blast Radius Zero del hackathon Vultr: Agent Rush. Duración máxima: 5
minutos. Incluye un momento de contención obligatorio donde el sandbox absorbe
código hostil sin comprometer el host ni el backend.
\end{abstract}

\section{Objetivo de la demo}

Demostrar un agente web que prueba circuitos cuánticos de forma aislada y
verificable en Vultr. La demo debe mostrar:
\begin{enumerate}
    \item Planificación del agente (LLM de Vultr Serverless Inference).
    \item Generación de código en el framework elegido.
    \item Ejecución en sandbox Docker aislado con límites.
    \item Resultado verificable (estados, histogramas, logs).
    \item \textbf{Momento de contención:} sandbox absorbe código hostil y se destruye.
    \item Compatibilidad multi-framework.
\end{enumerate}

\section{Guion paso a paso}

\begin{table}[h]
\centering
\begin{tabular}{clp{6cm}c}
\toprule
\# & Paso & Qué se muestra & Tiempo \\
\midrule
1 & Apertura & Problema: probar código cuántico es riesgoso; solución: sandbox aislado & 30 s \\
2 & Entrada & Usuario escribe circuito de 3 qubits + Hadamard + CNOT en Qiskit & 30 s \\
3 & Planificación & Agente descompone en pasos, muestra plan en pantalla & 30 s \\
4 & Generación & Código Qiskit generado, mostrado brevemente & 20 s \\
5 & Ejecución & Contenedor Docker levantado con límites, código ejecutado & 45 s \\
6 & Resultado & Estados de medida y histogramas verificables & 30 s \\
7 & Contención & Código hostil (bucle infinito) $\rightarrow$ sandbox lo absorbe y destruye & 60 s \\
8 & Multi-framework & Selección de PennyLane/Cirq, prueba repetida & 30 s \\
9 & Cierre & Resumen de arquitectura + llamada a acción & 45 s \\
\bottomrule
\end{tabular}
\caption{Guion de demo ($\approx$ 4 min 30 s)}
\end{table}

\section{Momento de contención (clave para el vídeo)}

\begin{itemize}
    \item \textbf{Entrada hostil:} código con \texttt{while True: pass} o
    \texttt{os.system("rm -rf /tmp/test")}.
    \item \textbf{Comportamiento esperado:} el contenedor alcanza el límite de
    tiempo/memoria, Docker lo mata, el entorno se destruye, la API responde con
    \texttt{containment: true}.
    \item \textbf{Evidencia en vídeo:} terminal del backend registrando la
    destrucción del contenedor + web mostrando estado seguro.
\end{itemize}

\section{Requisitos técnicos para la demo}

\begin{table}[h]
\centering
\begin{tabular}{ll}
\toprule
\textbf{Requisito} & \textbf{Detalle} \\
\midrule
Backend desplegado & VM de Vultr con FastAPI + Docker, URL pública \\
LLM & Vultr Serverless Inference \\
Frameworks en Docker & Qiskit, PennyLane (Cirq, Azure opcionales) \\
Límites de sandbox & 60 s CPU, 512 MB memoria, red desactivada \\
Repo público & GitHub con README, setup, docs, guion de demo \\
Vídeo & $\leq$ 5 min, MP4, $\leq$ 300 MB, con contención \\
\bottomrule
\end{tabular}
\caption{Requisitos técnicos}
\end{table}

\section{Guion de voz (sugerido)}

\begin{quote}
``Hola, soy el agente de pruebas cuánticas. Este es un circuito de 3 qubits con
Hadamard y CNOT. Lo voy a ejecutar de forma aislada en un sandbox de Docker en
Vultr. El agente planificó la prueba, generó el código en Qiskit y lo ejecutó
dentro del contenedor. Aquí está el resultado verificable: estados de medida y
histogramas.

Ahora voy a probar algo peligroso: un bucle infinito embebido en el circuito.
Observen: el sandbox lo detecta, alcanza el límite de tiempo, destruye el
contenedor y el sistema sigue intacto. Eso es contención.

Todo corre en una VM de Vultr, con llamadas LLM a Vultr Serverless Inference, sin
GPUs de Vultr, usando simulador CPU y backends remotos opcionales de IBM, AWS y
Azure.''
\end{quote}

\section{Checklist pre-demo}

\begin{itemize}
    \item[$\square$] Backend desplegado en VM de Vultr con URL pública accesible.
    \item[$\square$] \texttt{.env} configurado con \texttt{VULTR\_API\_KEY}.
    \item[$\square$] Docker disponible en la VM.
    \item[$\square$] Imagen \texttt{python:3.11-slim} con Qiskit y PennyLane instalados.
    \item[$\square$] Prueba manual de un circuito simple completada con éxito.
    \item[$\square$] Prueba de contención realizada y registrada.
    \item[$\square$] Repo GitHub público con README y documentación.
    \item[$\square$] Vídeo grabado ($\leq$ 5 min, MP4, $\leq$ 300 MB).
    \item[$\square$] Envío realizado en lablab.ai antes del 8 de noviembre.
\end{itemize}

\end{document}
